\section{Pico 소스 코드별 API Reference}

본 부록은 \code{Pico/} 디렉터리 MicroPython 소스의 클래스\textbf{·}메서드\textbf{·}
함수 시그니처를 파일별로 정리한다. 시그니처는 소스에서 그대로 옮긴 것이며,
\code{backup/} 및 날짜 백업 파일은 제외한다.

\subsection{main.py}

\noindent\textbf{모듈 상수}: \code{UART\_ID=0}, \code{UART\_BAUDRATE=9600},
\code{MAIN\_LOOP\_DELAY\_MS=10}, \code{RECEIVE\_TIMEOUT\_MS=500},
\code{MAX\_BUFFER\_SIZE=512}.

\noindent\textbf{class AresRover} --- UART 통신과 명령 처리를 총괄하는 메인 컨트롤러.

\begin{longtable}{S{7.4cm} L{6.6cm}}
\toprule
\rowcolor{tablehead}\normalfont\sffamily\bfseries 시그니처 & \normalfont\sffamily\bfseries 설명 \\
\midrule
\endhead
def \_\_init\_\_(self): & UART/명령처리기/하드웨어 초기화 \\
def boot(self): & 안전 정지\textbf{·}부팅음\textbf{·}OLED 메시지 부팅 시퀀스 \\
def \_safe\_stop\_all\_motors(self): & 전체 모터 비상 정지 \\
def \_pop\_line(self): & 버퍼에서 개행 단위 완성 라인 추출(없으면 None) \\
def \_read\_uart\_line(self): & 타임아웃 허용 다중 청크 UART 라인 수신 \\
def \_needs\_response(self, data): & 응답 필요 여부 판단(fire-and-forget 필터) \\
def \_is\_status\_message(self, data): & '+' 시작 BT 상태 메시지 여부 \\
def \_send\_response(self, response): & 20바이트 청크로 응답 전송 \\
def \_process\_command(self, data): & 명령 처리 후 결과 문자열 반환 \\
def run(self): & 명령 수신\textbf{·}부저 갱신 메인 루프 \\
\bottomrule
\end{longtable}

\subsection{process\_data.py}

\noindent\textbf{모듈 상수}: \code{PWM\_MAX=65535}.

\noindent\textbf{class CommandProcessor} --- UART 명령 파서 및 디스패처.

\begin{longtable}{S{7.8cm} L{6.2cm}}
\toprule
\rowcolor{tablehead}\normalfont\sffamily\bfseries 시그니처 & \normalfont\sffamily\bfseries 설명 \\
\midrule
\endhead
def process(self, data): & 명령을 핸들러로 분배하는 메인 디스패처 \\
def \_get\_speed\_pwm(self): & max\_speed로부터 PWM 값 계산 \\
def \_handle\_ping(self): & 연결 확인 + OLED 갱신 \\
def \_handle\_stop\_all(self): & 비상 정지 + OLED 메시지 \\
def \_handle\_get\_status(self): & 거리\textbf{·}자기\textbf{·}온도\textbf{·}메모리\textbf{·}업타임 반환 \\
def \_handle\_get\_sys(self): & 시스템 설정값 반환 \\
def \_handle\_get\_modules(self): & 활성 모듈 정보 반환 \\
def \_handle\_set\_pin(self, data): & 핀 할당 변경(SET\_PIN,이름,번호) \\
def \_handle\_set\_module(self, data): & 모듈 활성/비활성(SET\_MODULE,이름,0/1) \\
def \_handle\_batch(self, data): & 파이프 구분 배치 실행(BATCH;cmd1|cmd2) \\
def \_handle\_sys\_set(self, data): & 시스템 설정 저장(SYS\_SET,...) \\
def \_handle\_buzzer\_on(self, data): & 비차단 부저 시작(BUZZER\_ON,주파수,지속) \\
def \_handle\_sing(self): & 멜로디 재생(Hala Madrid) \\
def \_handle\_timed\_wheel(self, data, direction): & 방향별 시간 지정 서보 이동 \\
def \_handle\_continuous\_wheel(self, direction): & 연속 서보 이동 \\
def \_handle\_dcmotor(self, data): & 레거시 DC 모터 제어 \\
def \_handle\_timed\_dcmotor(self, data): & 레거시 시간 지정 DC 모터 \\
def \_handle\_timed\_dcmotor\_new(self, data, direction): & 신규 형식 시간 지정 DC 모터 \\
def \_handle\_main\_motor(self, direction): & 연속 DC 모터 구동 \\
def \_handle\_calib\_start(self): & 보정 테스트 시퀀스 실행 \\
def \_handle\_calib\_set(self, data): & 보정값 저장(CALIB\_SET,좌,우) \\
def \_handle\_led\_on(self, data): & LED 켜기(LED\_ON,번호,밝기) \\
def \_handle\_led\_off(self, data): & LED 끄기(LED\_OFF,번호|ALL) \\
def \_handle\_led\_pattern(self, data): & LED 패턴 적용([v0 v1 ... v5]) \\
def \_handle\_msg(self, data): & OLED 텍스트(16자 자동 줄바꿈) \\
def \_handle\_clear\_rect(self, data): & 사각 영역 지우기(CLEAR\_RECT,x,y,w,h) \\
def \_handle\_msg\_xy(self, data): & 좌표 텍스트(MSG\_XY,x,y,text) \\
def \_handle\_icon(self, data): & 아이콘 표시(ICON,name,x,y) \\
def \_handle\_sleep(self, data): & 대기(SLEEP,초) \\
\bottomrule
\end{longtable}

\subsection{hardware.py}

\noindent\textbf{class RobotHardware} --- 모듈 선택 로딩을 수행하는 하드웨어 싱글턴.
전역 인스턴스 \code{robot = RobotHardware()}.

\begin{longtable}{S{7.4cm} L{6.6cm}}
\toprule
\rowcolor{tablehead}\normalfont\sffamily\bfseries 시그니처 & \normalfont\sffamily\bfseries 설명 \\
\midrule
\endhead
def \_\_new\_\_(cls): & 싱글턴 인스턴스 보장 \\
def \_init\_hardware(self): & 설정에 따른 전체 모듈 초기화 \\
def \_sync\_bt\_name(self): & 설정의 블루투스 이름 동기화 \\
def get\_temperature(self): & 내장 온도 센서 읽기 \\
def get\_battery\_voltage(self): & 배터리 전압(미구현, 100 반환) \\
def get\_status\_dict(self): & 센서값\textbf{·}지표 상태 딕셔너리 반환 \\
def stop\_all(self): & 모터\textbf{·}LED\textbf{·}부저\textbf{·}발사기 전체 정지 \\
\bottomrule
\end{longtable}

\subsection{system\_config.py}

\noindent\textbf{모듈 상수}: \code{CONFIG\_FILE='system.json'},
\code{LEGACY\_CALIB\_FILE='calibration.json'},
\code{MAX\_DEVICE\_NAME\_LENGTH=10}, \code{DEFAULT\_CONFIG}, \code{CONFIG\_KEY\_ORDER}.

\noindent\textbf{class SystemConfig} --- 파일 영속화 설정 싱글턴.
전역 인스턴스 \code{sys\_config = SystemConfig()}.

\begin{longtable}{S{7.8cm} L{6.2cm}}
\toprule
\rowcolor{tablehead}\normalfont\sffamily\bfseries 시그니처 & \normalfont\sffamily\bfseries 설명 \\
\midrule
\endhead
def \_\_new\_\_(cls): & 싱글턴 인스턴스 \\
def \_init\_config(self): & 설정 초기화 및 파일 로드 \\
def load\_config(self): & system.json 로드(레거시 보정 파일 이관 포함) \\
def save\_config(self, max\_speed, collision\_dist, auto\_stop, device\_name): & 시스템 설정 저장 \\
def save\_calibration(self, left, right): & 휠 보정 계수 저장 \\
def save\_all(self): & 순서화 키로 전체 설정 저장 \\
def get(self, key): & 설정값 조회 \\
def set\_pin(self, pin\_name, pin\_number): & 핀 할당 변경(0\textasciitilde 28 검증) \\
def enable\_module(self, module\_name, enable): & 모듈 활성/비활성 \\
def is\_module\_enabled(self, module\_name): & 모듈 활성 여부 \\
def get\_module\_info(self): & 모듈 상태\textbf{·}설명 반환 \\
\bottomrule
\end{longtable}

\subsection{pins.py}

\noindent\textbf{모듈 상수}:
\begin{lstlisting}[style=json]
UART_TX_PIN=0   UART_RX_PIN=1     DCMOTOR_DIR_PIN=8   DCMOTOR_PWM_PIN=9
I2C_SDA_PIN=4   I2C_SCL_PIN=5     ULTRASONIC_TRIG_PIN=6  ULTRASONIC_ECHO_PIN=7
SERVO_RIGHT_PIN=12   SERVO_LEFT_PIN=13   BUZZER_PIN=15   GUN_PIN=22
MAGSENSOR_PIN=26   LED_PINS=[21,20,19,18,17,16]   RESERVED_PINS=[10,11,14,27,28]
\end{lstlisting}

\noindent\textbf{함수}: \code{def print\_pin\_info():} --- 핀 할당 표를 콘솔 출력.

\subsection{wheel.py}

\noindent\textbf{class KSWheel} --- 보정 지원 서보 휠 컨트롤러.

\begin{longtable}{S{7.4cm} L{6.6cm}}
\toprule
\rowcolor{tablehead}\normalfont\sffamily\bfseries 시그니처 & \normalfont\sffamily\bfseries 설명 \\
\midrule
\endhead
def \_\_init\_\_(self): & 서보 중립 듀티\textbf{·}보정 계수 초기화 \\
def update\_factors(self, left, right): & 보정 계수 갱신(최대 100) \\
def set\_angle(self, wheel\_idx, angle, speed=1.0): & 각도(0\textasciitilde 180)→PWM 듀티 매핑 \\
def forward(self, speed=0.1): & 전진 \\
def backward(self, speed=0.1): & 후진 \\
def turn\_right(self, speed=0.1): & 우회전 \\
def turn\_left(self, speed=0.1): & 좌회전 \\
def stop(self): & 양 휠 정지(중립) \\
\bottomrule
\end{longtable}

\subsection{dcmotor.py}

\noindent\textbf{모듈 상수}: \code{PWM\_FREQUENCY=20000}, \code{PWM\_STOP=0},
\code{PWM\_MAX=65535}, \code{DEBUG\_MOTOR=False}.

\noindent\textbf{class KSDCMotor} --- 브레이크 변조 듀얼 PWM H-브리지 DC 모터.

\begin{longtable}{S{7.4cm} L{6.6cm}}
\toprule
\rowcolor{tablehead}\normalfont\sffamily\bfseries 시그니처 & \normalfont\sffamily\bfseries 설명 \\
\midrule
\endhead
def \_\_init\_\_(self): & IN1\textbf{·}IN2 양 핀 PWM 초기화 \\
def dc\_forward(self, speed): & 전진(IN2를 PWM\_MAX-speed로 변조) \\
def dc\_backward(self, speed): & 후진(역방향 변조) \\
def dc\_stop(self): & 정지(코스트) \\
def is\_running(self): & 구동 상태 반환 \\
\bottomrule
\end{longtable}

\subsection{buzzer.py}

\noindent\textbf{class KSBuzzer} --- 차단/비차단 재생 부저.

\begin{longtable}{S{7.6cm} L{6.4cm}}
\toprule
\rowcolor{tablehead}\normalfont\sffamily\bfseries 시그니처 & \normalfont\sffamily\bfseries 설명 \\
\midrule
\endhead
def \_\_init\_\_(self, pin\_num=None): & 부저 PWM 초기화(기본 BUZZER\_PIN) \\
def play(self, freq=1000, duration=1, vol=1000): & 완료까지 차단 재생 \\
def start(self, freq=1000, duration=1, vol=1000): & 비차단 시작(재생 중이면 False) \\
def update(self): & 지속시간 경과 시 자동 정지 \\
def stop(self): & 즉시 정지 \\
def boot\_sound(self, short=False): & 부팅음(G5 0.2s + C4 0.9s) \\
def halamadrid(self, short=False): & 18음 멜로디 재생 \\
\bottomrule
\end{longtable}

\subsection{leds.py}

\noindent\textbf{class KSLeds} --- 패턴 지원 LED 배열 컨트롤러.

\begin{longtable}{S{7.4cm} L{6.6cm}}
\toprule
\rowcolor{tablehead}\normalfont\sffamily\bfseries 시그니처 & \normalfont\sffamily\bfseries 설명 \\
\midrule
\endhead
def \_\_init\_\_(self): & 6개 LED 20kHz PWM 초기화 후 소등 \\
def leds\_off(self): & 전체 소등 \\
def swipe\_effect(self): & 스와이프 애니메이션(100ms 타이밍) \\
def check(self): & 시작 점검 패턴(3패스, 페이딩) \\
def set\_led\_pattern(self, pattern): & 밝기 패턴 적용(0.0\textasciitilde 1.0 리스트) \\
\bottomrule
\end{longtable}

\subsection{gun.py}

\noindent\textbf{class KSGun} --- BB탄 발사기 컨트롤러.

\begin{longtable}{S{7.8cm} L{6.2cm}}
\toprule
\rowcolor{tablehead}\normalfont\sffamily\bfseries 시그니처 & \normalfont\sffamily\bfseries 설명 \\
\midrule
\endhead
def \_\_init\_\_(self, pin\_num=None): & 발사 모터 PWM 1kHz 초기화(기본 GUN\_PIN) \\
def power(self, value): & PWM 듀티 설정 \\
def stop(self): & 모터 정지 \\
def soft\_start(self, target\_power, duration\_ms=150): & 8단계 파워 램프업 \\
def fire\_once(self, power=50000, spin\_time=0.22, cooldown\_ms=250): & 단발 발사(램프업→회전→정지) \\
\bottomrule
\end{longtable}

\subsection{ultrasonic.py}

\noindent\textbf{class KSDistance} --- HC-SR04 초음파 거리 센서.

\begin{longtable}{S{7.4cm} L{6.6cm}}
\toprule
\rowcolor{tablehead}\normalfont\sffamily\bfseries 시그니처 & \normalfont\sffamily\bfseries 설명 \\
\midrule
\endhead
def \_\_init\_\_(self): & 트리거(OUT)\textbf{·}에코(IN) 핀 초기화 \\
def get\_distance(self): & cm 거리 측정(30ms 타임아웃, 오류 시 1000000) \\
\bottomrule
\end{longtable}

\subsection{magsensor.py}

\noindent\textbf{class KSMagSensor} --- 자기장 센서(액티브 로우).

\begin{longtable}{S{7.4cm} L{6.6cm}}
\toprule
\rowcolor{tablehead}\normalfont\sffamily\bfseries 시그니처 & \normalfont\sffamily\bfseries 설명 \\
\midrule
\endhead
def \_\_init\_\_(self): & MAGSENSOR\_PIN 풀업 입력 초기화 \\
def detect(self): & 자석 검출(0→검출=1, 1→미검출=0 반환) \\
\bottomrule
\end{longtable}

\subsection{ssd1306.py}

\noindent SSD1306 명령 레지스터 상수(\code{SET\_CONTRAST=0x81} 등) 다수 정의.

\noindent\textbf{class SSD1306} --- FrameBuffer 기반 OLED 베이스 컨트롤러.

\begin{longtable}{S{7.6cm} L{6.4cm}}
\toprule
\rowcolor{tablehead}\normalfont\sffamily\bfseries 시그니처 & \normalfont\sffamily\bfseries 설명 \\
\midrule
\endhead
def \_\_init\_\_(self, width, height, external\_vcc): & 프레임버퍼\textbf{·}아이콘 초기화 \\
def init\_display(self): & 초기화 명령 시퀀스 전송 \\
def poweroff(self): / def poweron(self): & 디스플레이 전원 제어 \\
def contrast(self, contrast): & 명암 설정(0\textasciitilde 255) \\
def invert(self, invert): & 색 반전 \\
def show(self): & 프레임버퍼 플러시 \\
def fill(self, col): & 전체 채우기(0/1) \\
def pixel(self, x, y, col): & 픽셀 설정 \\
def fill\_rect(self, x, y, w, h, col): & 사각 영역 채우기 \\
def scroll(self, dx, dy): & 프레임버퍼 스크롤 \\
def text(self, string, x, y, col=1): & 텍스트 그리기 \\
def booting\_msg(self): & 부팅 스플래시 표시 \\
\bottomrule
\end{longtable}

\noindent\textbf{class SSD1306\_I2C} --- I2C 인터페이스 변형.

\begin{longtable}{S{9.6cm} L{4.4cm}}
\toprule
\rowcolor{tablehead}\normalfont\sffamily\bfseries 시그니처 & \normalfont\sffamily\bfseries 설명 \\
\midrule
\endhead
def \_\_init\_\_(self, width, height, scl\_pin=None, sda\_pin=None, addr=0x3c, external\_vcc=False): & I2C 초기화 \\
def write\_cmd(self, cmd): & I2C 명령 바이트 쓰기 \\
def write\_data(self, buf): & I2C 데이터 버퍼 쓰기 \\
\bottomrule
\end{longtable}

\noindent\textbf{class SSD1306\_SPI} --- SPI 인터페이스 변형(현재 미사용).
\code{def \_\_init\_\_(self, width, height, spi, dc, res, cs, external\_vcc=False):},
\code{write\_cmd(self, cmd)}, \code{write\_data(self, buf)}.

\subsection{icon.py}

\noindent\textbf{모듈 상수(비트맵)}: \code{mars\_rover32x32}, \code{cute\_robot32x32},
\code{open\_eye32x32}, \code{closed\_eye32x32} (각 32$\times$32).

\noindent\textbf{class KSicon} --- 수평 바이트 비트맵을 MONO\_VLSB 수직 바이트로 변환\textbf{·}렌더.

\begin{longtable}{S{7.4cm} L{6.6cm}}
\toprule
\rowcolor{tablehead}\normalfont\sffamily\bfseries 시그니처 & \normalfont\sffamily\bfseries 설명 \\
\midrule
\endhead
def \_\_init\_\_(self, buff, w, h, oled): & 비트맵 포맷 변환 \\
def blit(self, x, y): & 지정 좌표에 아이콘 렌더 \\
\bottomrule
\end{longtable}
